\section*{Ensayo por Andrea Lucia Olmos Ortiz}
\addcontentsline{toc}{subsubsection}{Ensayo por Andrea Lucia Olmos Ortiz.}
\subsection*{Introducción.}
Con el inicio de una nueva revolución industrial, marcada por la economía digital, en la que la información ha adquirido un valor monetario, debemos preguntarnos: ¿cuánto vale nuestra información personal?\\ \\
Aunque actualmente me considero consciente de la información que comparto y de sus implicaciones, he estado activa en redes sociales incluso antes de la publicación de este artículo y puedo decir con toda seguridad que en su momento no fui consciente de la cantidad de datos personales que compartía, ni del uso que probablemente se les dio. En este contexto, me parece esencial que las personas conozcan y reflexionen sobre el contenido de este artículo.
\subsection*{Desarrollo.}
La autora afirma que la esencia de todas las redes sociales es intercambiar nuestras ideas, experiencias y hábitos, como información que permite a las empresas mejorar sus estrategias de publicidad, a cambio de no sentirnos excluidos. Estoy completamente de acuerdo con esta afirmación, y es algo que he podido observar día a día durante toda mi adolescencia, donde gran parte de la convivencia está basada en lo que sucede en las redes sociales. Conforme más lo cuestionaba, más me daba cuenta de lo injusto de esa “transacción", ya que pertenecer es una de las necesidades más humanas que tenemos.\\ \\
Por otra parte, el artículo presenta los entornos de \textit{Internet of Things} como uno de los principales factores detrás del crecimiento exponencial de los datos generados diariamente, lo cual exige un manejo adecuado Big Data para poder aprovecharlos de forma efectiva. \\ \\
También resalta que la economía digital continua creciendo gracias a la información que nosotros mismos generamos como usuarios, la cual también crece exponencialmente. Esta parte del artículo me resonó con los conceptos que revisamos en las primeras clases de la materia, la enorme cantidad de datos que se generan cada año y la importancia de conocer métodos adecuados para organizarlos, almacenarlos y aprovecharlos. \\ \\
Por último, el artículo intenta responder la pregunta central acerca del valor de la información desde dos enfoques.\\ \\
1. El valor que las compañías le dan a nuestros datos, el cual puede medirse con el ingreso promedio por usuario(ARPU); por ejemplo, para Google fue de aproximadamente 40 USD por usuario.\\ \\
2. La percepción individual del valor de nuestra información.\\ \\
En este segundo enfoque, la información se vuelve un poco caótica, ya que proviene de múltiples estudios y estos no parecen presentar resultados consistentes. \\ \\
\subsection*{Conclusión.}
Llegué a la conclusión de que independientemente del valor intrínseco que las compañias le den a nuestros datos la percepción individual del valor de nuestra información depende de la importancia que nosotros mismos le demos. Esto reafirma la relevancia de comprender las implicaciones económicas de los datos que compartimos y de ser conscientes de qué información estamos dispuestos a dar y cuál preferimos mantener privada.\\ \\
Me pareció un artículo muy valioso para su tiempo, pero con ya casi diez años desde su publicación, es inevitable sentir curiosidad por las nuevas implicaciones que tiene nuestra información frente a los avances tecnológicos recientes, como la inteligencia artificial. Esto me lleva a cuestionarme si realmente sigue siendo suficiente con estar conscientes del uso de nuestra información, o si tal vez deberíamos estar haciendo algo más para proteger nuestros datos.